\section{Proposed research directions}
\label{sec:directions}

Based on the research gaps identified in Section~\ref{sec:gaps}, we propose three actionable research directions tailored to software engineering practice and academic capstone environments, complemented by a preliminary pilot validation:

\subsection{Empirical evaluation of enriched user stories using standardized competence rubrics}
While Keshk et al. \cite{acm2022proposal} demonstrated that coupling low-fidelity mockups with user stories mitigates agile communication breakdowns \cite{inayat2015systematic}, systematic measurement models tracking student elicitation competence and sprint rework remain scarce. Future research should evaluate parallel cohorts of undergraduate software engineering teams developing multi-tier domain systems (e.g., self-storage facility management or booking platforms). 

By adopting Lending et al.'s \cite{lending2022rubric} validated 5-level, 8-skill elicitation rubric, researchers can objectively measure student progression from Level 2 (passive prototype presentations) to Level 4 (collaborative interactive co-design). The primary dependent variables should include: (1) sprint velocity, (2) formal story clarification requests during development, (3) post-sprint defect rework effort, and (4) longitudinal Ambiguity Resolution Rate ($ARR$) across agile release cycles.

\subsection{Lightweight ambiguity checklists and automated LLM traceability pipelines}
While NLP ambiguity tools provide classification capabilities \cite{cascon2025ambiguity}, practitioners require actionable heuristics to overcome the cognitive trap of unconscious disambiguation \cite{berry2003contract}. A promising research direction is the development and validation of a two-stage hybrid backlog grooming workflow. In the manual stage, development teams rapidly cross-check user story statements against interactive prototype screens during sprint backlog grooming, specifically verifying that pronouns reference concrete visual elements and numerical boundary conditions are verifiable. In the automated stage, building upon Kolthoff et al. \cite{kolthoff2024interlinking}, LLMs automatically cross-reference acceptance criteria in user stories against Figma component trees and GUI mockups synthesized by generative tools like UI-Diffuser \cite{wei2023boosting} or SERGUI \cite{kolthoff2024selfelicitation}. This establishes bidirectional traceability and automatically alerts teams to unrepresented business rules or ambiguous terminology before implementation begins.

\subsection{Augmenting wireframes with think-aloud and in-situ protocols for quality attributes}
Because static visual mockups trigger an acute non-functional requirement deficit \cite{rueda2020requirements, davis2018maturity}, future empirical protocols must combine cognitive and operational contextual augmentations. Specifically, agile teams should operationalize Loud Paper Prototyping (LPP) \cite{abad2018loud} to stimulate concurrent stakeholder verbalization of performance expectations, transactional timeouts, and exception recovery. 

Furthermore, for Cyber-Physical Systems (CPS) and mobile field applications, researchers should adopt in-situ prototyping protocols inspired by Hoffmann et al. \cite{hoffmann2011automotive}. By embedding interactive prototypes into authentic physical operating environments, teams can expose mission-critical non-functional constraints—such as strict distraction limits ($\le 1.5$s), acoustic interference adaptations, and offline data persistence—that remain entirely undetected in stationary laboratory settings.

\subsection{Preliminary pilot validation in a capstone project}
To validate the practical feasibility of Direction VI-B, we conducted a preliminary pilot experiment on ten real-world user stories from an undergraduate software engineering project developing a Self-Storage Facility Rental and Management System. Requirements ambiguity was evaluated across the four linguistic categories established by Berry et al. \cite{berry2003contract}. In the baseline textual user stories, fifteen distinct ambiguity defects were identified across all ten statements (100\% defect rate). 

Following the introduction of low-fidelity interactive wireframes and visual acceptance criteria, fourteen ambiguities were successfully eliminated, achieving an Ambiguity Resolution Rate of $ARR = 93.3\%$ (structural, scope, and anaphoric ambiguities were reduced by 100\%). A non-parametric Wilcoxon signed-rank test confirmed a statistically significant reduction in requirement defects ($Z = -2.81, p = 0.005 < 0.01$). Crucially, the single unresolved defect involved transaction latency, empirically corroborating the non-functional requirement deficit of static visual prototypes \cite{rueda2020requirements} and highlighting the necessity of combining wireframes with explicit quality attribute checklists.
